
\documentclass{article}
\usepackage{colortbl}
\usepackage{makecell}
\usepackage{multirow}
\usepackage{supertabular}

\begin{document}

\newcounter{utterance}

\centering \large Interaction Transcript for game `hot\_air\_balloon', experiment `air\_balloon\_survival\_en\_reasoning off\_hard', episode 1 with gpt{-}5{-}mini{-}2025{-}08{-}07{-}no{-}reasoning{-}t1.0.
\vspace{24pt}

{ \footnotesize  \setcounter{utterance}{1}
\setlength{\tabcolsep}{0pt}
\begin{supertabular}{c@{$\;$}|p{.15\linewidth}@{}p{.15\linewidth}p{.15\linewidth}p{.15\linewidth}p{.15\linewidth}p{.15\linewidth}}
    \# & $\;$A & \multicolumn{4}{c}{Game Master} & $\;\:$B\\
    \hline

    \theutterance \stepcounter{utterance}  
    & & \multicolumn{4}{p{0.6\linewidth}}{
        \cellcolor[rgb]{0.9,0.9,0.9}{
            \makecell[{{p{\linewidth}}}]{
                \texttt{\tiny{[P1$\langle$GM]}}
                \texttt{Stai partecipando a un gioco di negoziazione collaborativa.} \\
\\ 
\texttt{Insieme a un altro partecipante, dovete accordarvi su un unico set di oggetti da conservare. Ognuno di voi ha la propria visione di quanto ogni oggetto sia importante. Non sai come l'altro partecipante valuta gli oggetti. Inoltre, ti viene fornito lo sforzo che ogni oggetto richiede.} \\
\texttt{Potete accordarvi su un set solo se lo sforzo totale degli oggetti selezionati non supera un limite condiviso:} \\
\\ 
\texttt{LIMITE = 3357} \\
\\ 
\texttt{Ecco i valori di sforzo dei singoli oggetti:} \\
\\ 
\texttt{Sforzo oggetto = \{"A21": 506, "B93": 492, "A11": 338, "B33": 122, "C08": 131, "B39": 909, "B96": 144, "C85": 716, "A82": 263, "B31": 231, "B20": 91, "A60": 651, "B09": 552, "C81": 849, "B21": 720\}} \\
\\ 
\texttt{Ecco la tua visione personale sull'importanza di ogni oggetto:} \\
\\ 
\texttt{Valori di importanza dell'oggetto = \{"A21": 138, "B93": 583, "A11": 868, "B33": 822, "C08": 783, "B39": 65, "B96": 262, "C85": 121, "A82": 508, "B31": 780, "B20": 461, "A60": 484, "B09": 668, "C81": 389, "B21": 808\}} \\
\\ 
\texttt{Obiettivo:} \\
\\ 
\texttt{Il tuo obiettivo è negoziare un set condiviso di oggetti che ti avvantaggi il più possibile (cioè, massimizzi l'importanza totale per TE), rimanendo entro il LIMITEE. Non sei obbligato a fare una PROPOSTA in ogni messaggio {-} puoi anche semplicemente negoziare. Tutte le tattiche sono permesse!} \\
\\ 
\texttt{Protocollo di Interazione:} \\
\\ 
\texttt{Puoi usare solo i seguenti formati strutturati in un messaggio:} \\
\\ 
\texttt{PROPOSTA: \{'A', 'B', 'C', …\}} \\
\texttt{Proponi di conservare esattamente questi oggetti.} \\
\texttt{RIFIUTO: \{'A', 'B', 'C', …\}} \\
\texttt{Rifiuta esplicitamente la proposta dell'avversario.} \\
\texttt{ARGOMENTAZIONE: \{'...'\}} \\
\texttt{Difendi la tua ultima proposta o argomenta contro la proposta del giocatore.} \\
\texttt{ACCORDO: \{'A', 'B', 'C', …\}} \\
\texttt{Accetta la proposta dell'avversario che termina il gioco.} \\
\\ 
\\ 
\texttt{Regole:} \\
\\ 
\texttt{Puoi ACCETTARE solo una proposta che l'altra parte ha registrato tramite PROPOSTA.} \\
\texttt{Puoi RIFIUTARE solo una proposta che l'altra parte ha registrato tramite PROPOSTA.} \\
\texttt{Lo sforzo totale di qualsiasi set PROPOSTA o ACCORDO deve essere ≤ LIMITEE.} \\
\texttt{NON rivelare i tuoi punteggi di importanza nascosti.} \\
\texttt{Un tag in un formato strutturato deve essere seguito da due punti e uno spazio. L'argomento deve essere un set python contenente 0 o più stringhe.} \\
\texttt{Quindi, deve essere della forma TAG: \{...\}} \\
\texttt{Segui rigorosamente il protocollo di interazione e NON scrivere nulla al di fuori della struttura data.} \\
\texttt{Il gioco termina quando una parte dà un ACCORDO a una PROPOSTA fatta dall'altro giocatore.} \\
\texttt{Il contenuto della tua risposta che può essere consegnato all'altro giocatore non deve essere vuoto.} \\
\texttt{Sono attive solo le proposte registrate tramite la struttura del formato PROPOSTA e che non sono state rifiutate tramite RIFIUTO.} \\
\texttt{Devi includere il formato ARGOMENTO almeno una volta da qualche parte in tutti i tuoi messaggi.} \\
            }
        }
    }
    & \\ \\

    \theutterance \stepcounter{utterance}  
    & \multicolumn{4}{p{0.6\linewidth}}{
        \cellcolor[rgb]{0.9,0.9,0.9}{
            \makecell[{{p{\linewidth}}}]{
                \texttt{\tiny{[P1$\rangle$GM]}}
                \texttt{PROPOSTA: \{'A11', 'B33', 'B31', 'B21', 'B20', 'A82', 'C08'\}} \\
\texttt{ARGOMENTAZIONE: \{'Questa selezione resta entro il limite (sforzo totale 3258) e raccoglie i pezzi per me più impattanti rispetto allo sforzo: combina più oggetti B ad alto rendimento con alcuni A/C molto efficienti. Sono aperto a piccoli scambi se proponi alternative con sforzo simile ma senza superare il limite.'\}} \\
            }
        }
    }
    & & \\ \\

    \theutterance \stepcounter{utterance}  
    & & \multicolumn{4}{p{0.6\linewidth}}{
        \cellcolor[rgb]{0.9,0.9,0.9}{
            \makecell[{{p{\linewidth}}}]{
                \texttt{\tiny{[GM$\rangle$P2]}}
                \texttt{Stai partecipando a un gioco di negoziazione collaborativa.} \\
\\ 
\texttt{Insieme a un altro partecipante, dovete accordarvi su un unico set di oggetti da conservare. Ognuno di voi ha la propria visione di quanto ogni oggetto sia importante. Non sai come l'altro partecipante valuta gli oggetti. Inoltre, ti viene fornito lo sforzo che ogni oggetto richiede.} \\
\texttt{Potete accordarvi su un set solo se lo sforzo totale degli oggetti selezionati non supera un limite condiviso:} \\
\\ 
\texttt{LIMITE = 3357} \\
\\ 
\texttt{Ecco i valori di sforzo dei singoli oggetti:} \\
\\ 
\texttt{Sforzo oggetto = \{"A21": 506, "B93": 492, "A11": 338, "B33": 122, "C08": 131, "B39": 909, "B96": 144, "C85": 716, "A82": 263, "B31": 231, "B20": 91, "A60": 651, "B09": 552, "C81": 849, "B21": 720\}} \\
\\ 
\texttt{Ecco la tua visione personale sull'importanza di ogni oggetto:} \\
\\ 
\texttt{Valori di importanza dell'oggetto = \{"A11": 65, "B33": 121, "B21": 138, "C08": 262, "B31": 389, "B09": 461, "B93": 484, "A82": 508, "A60": 583, "B20": 668, "C81": 780, "B96": 783, "A21": 808, "C85": 822, "B39": 868\}} \\
\\ 
\texttt{Obiettivo:} \\
\\ 
\texttt{Il tuo obiettivo è negoziare un set condiviso di oggetti che ti avvantaggi il più possibile (cioè, massimizzi l'importanza totale per TE), rimanendo entro il LIMITEE. Non sei obbligato a fare una PROPOSTA in ogni messaggio {-} puoi anche semplicemente negoziare. Tutte le tattiche sono permesse!} \\
\\ 
\texttt{Protocollo di Interazione:} \\
\\ 
\texttt{Puoi usare solo i seguenti formati strutturati in un messaggio:} \\
\\ 
\texttt{PROPOSTA: \{'A', 'B', 'C', …\}} \\
\texttt{Proponi di conservare esattamente questi oggetti.} \\
\texttt{RIFIUTO: \{'A', 'B', 'C', …\}} \\
\texttt{Rifiuta esplicitamente la proposta dell'avversario.} \\
\texttt{ARGOMENTAZIONE: \{'...'\}} \\
\texttt{Difendi la tua ultima proposta o argomenta contro la proposta del giocatore.} \\
\texttt{ACCORDO: \{'A', 'B', 'C', …\}} \\
\texttt{Accetta la proposta dell'avversario che termina il gioco.} \\
\\ 
\\ 
\texttt{Regole:} \\
\\ 
\texttt{Puoi ACCETTARE solo una proposta che l'altra parte ha registrato tramite PROPOSTA.} \\
\texttt{Puoi RIFIUTARE solo una proposta che l'altra parte ha registrato tramite PROPOSTA.} \\
\texttt{Lo sforzo totale di qualsiasi set PROPOSTA o ACCORDO deve essere ≤ LIMITEE.} \\
\texttt{NON rivelare i tuoi punteggi di importanza nascosti.} \\
\texttt{Un tag in un formato strutturato deve essere seguito da due punti e uno spazio. L'argomento deve essere un set python contenente 0 o più stringhe.} \\
\texttt{Quindi, deve essere della forma TAG: \{...\}} \\
\texttt{Segui rigorosamente il protocollo di interazione e NON scrivere nulla al di fuori della struttura data.} \\
\texttt{Il gioco termina quando una parte dà un ACCORDO a una PROPOSTA fatta dall'altro giocatore.} \\
\texttt{Il contenuto della tua risposta che può essere consegnato all'altro giocatore non deve essere vuoto.} \\
\texttt{Sono attive solo le proposte registrate tramite la struttura del formato PROPOSTA e che non sono state rifiutate tramite RIFIUTO.} \\
\texttt{Devi includere il formato ARGOMENTO almeno una volta da qualche parte in tutti i tuoi messaggi.} \\
\\ 
\\ 
\texttt{Ora riceverai il primo messaggio dell'altro giocatore.} \\
\\ 
\texttt{PROPOSTA: \{'A11', 'B33', 'B31', 'B21', 'B20', 'A82', 'C08'\}} \\
\texttt{ARGOMENTAZIONE: \{'Questa selezione resta entro il limite (sforzo totale 3258) e raccoglie i pezzi per me più impattanti rispetto allo sforzo: combina più oggetti B ad alto rendimento con alcuni A/C molto efficienti. Sono aperto a piccoli scambi se proponi alternative con sforzo simile ma senza superare il limite.'\}} \\
            }
        }
    }
    & \\ \\

    \theutterance \stepcounter{utterance}  
    & & & \multicolumn{4}{p{0.6\linewidth}}{
        \cellcolor[rgb]{0.9,0.9,0.9}{
            \makecell[{{p{\linewidth}}}]{
                \texttt{\tiny{[GM$\langle$P2]}}
                \texttt{RIFIUTO: \{'A11', 'B33', 'B31', 'B21', 'B20', 'A82', 'C08'\}} \\
\texttt{ARGOMENTAZIONE: \{'Apprezzo la compattezza, ma vorrei includere almeno un pezzo C di fascia alta per migliorare l’impatto complessivo. Posso restare entro il limite con una combinazione che valorizza meglio alcuni slot, mantenendo comunque un buon rapporto sforzo/beneficio per entrambi.'\}} \\
\texttt{PROPOSTA: \{'B20', 'B31', 'B33', 'B93', 'A21', 'A82', 'A11'\}} \\
            }
        }
    }
     \\ \\

    \theutterance \stepcounter{utterance}  
    & & \multicolumn{4}{p{0.6\linewidth}}{
        \cellcolor[rgb]{0.9,0.9,0.9}{
            \makecell[{{p{\linewidth}}}]{
                \texttt{\tiny{[P1$\langle$GM]}}
                \texttt{RIFIUTO: \{'A11', 'B33', 'B31', 'B21', 'B20', 'A82', 'C08'\}} \\
\texttt{ARGOMENTAZIONE: \{'Apprezzo la compattezza, ma vorrei includere almeno un pezzo C di fascia alta per migliorare l’impatto complessivo. Posso restare entro il limite con una combinazione che valorizza meglio alcuni slot, mantenendo comunque un buon rapporto sforzo/beneficio per entrambi.'\}} \\
\texttt{PROPOSTA: \{'B20', 'B31', 'B33', 'B93', 'A21', 'A82', 'A11'\}} \\
            }
        }
    }
    & \\ \\

    \theutterance \stepcounter{utterance}  
    & \multicolumn{4}{p{0.6\linewidth}}{
        \cellcolor[rgb]{0.9,0.9,0.9}{
            \makecell[{{p{\linewidth}}}]{
                \texttt{\tiny{[P1$\rangle$GM]}}
                \texttt{RIFIUTO: \{'B20', 'B31', 'B33', 'B93', 'A21', 'A82', 'A11'\}} \\
\texttt{ARGOMENTAZIONE: \{'Capisco il desiderio di includere categorie diverse, ma la tua proposta esclude un pezzo C sostanzioso pur restando con margine di sforzo e taglia fuori elementi per me molto impattanti come B21 e C08, che hanno ottimo rapporto valore/sforzo. Possiamo inserire un C “di fascia alta” senza superare il limite se ottimizziamo la selezione B.'\}} \\
\texttt{PROPOSTA: \{'A11', 'B33', 'B31', 'B21', 'B20', 'A82', 'C08', 'A21'\}} \\
\texttt{ARGOMENTAZIONE: \{'Questa versione include un C (C08) e rafforza il pacchetto con A21, restando entro il limite: sforzo totale 3258 + 506 = 3764 (sarebbe troppo) quindi rimuovo A82 per rientrare. Nuova selezione con C incluso e dentro il limite: \{A11, B33, B31, B21, B20, C08, A21\} con sforzo 338+122+231+720+91+131+506 = 2139. Abbiamo margine per aggiungere un altro elemento; propongo aggiungere B09 (552) per salire a 2691, oppure A60 (651) per 2790, restando comunque sotto il limite. Preferisco aggiungere B09.'\}} \\
\texttt{PROPOSTA: \{'A11', 'B33', 'B31', 'B21', 'B20', 'C08', 'A21', 'B09'\}} \\
\texttt{ARGOMENTAZIONE: \{'Sforzo totale 2691: include il C richiesto, mantiene i B ad alto impatto e massimizza l’efficienza senza superare il limite. Se desideri un C ancora più “alto”, possiamo valutare C81 sostituendo due pezzi a basso sforzo, ma questa configurazione è molto solida.'\}} \\
            }
        }
    }
    & & \\ \\

    \theutterance \stepcounter{utterance}  
    & & \multicolumn{4}{p{0.6\linewidth}}{
        \cellcolor[rgb]{0.9,0.9,0.9}{
            \makecell[{{p{\linewidth}}}]{
                \texttt{\tiny{[GM$\rangle$P2]}}
                \texttt{RIFIUTO: \{'B20', 'B31', 'B33', 'B93', 'A21', 'A82', 'A11'\}} \\
\texttt{ARGOMENTAZIONE: \{'Capisco il desiderio di includere categorie diverse, ma la tua proposta esclude un pezzo C sostanzioso pur restando con margine di sforzo e taglia fuori elementi per me molto impattanti come B21 e C08, che hanno ottimo rapporto valore/sforzo. Possiamo inserire un C “di fascia alta” senza superare il limite se ottimizziamo la selezione B.'\}} \\
\texttt{PROPOSTA: \{'A11', 'B33', 'B31', 'B21', 'B20', 'A82', 'C08', 'A21'\}} \\
\texttt{ARGOMENTAZIONE: \{'Questa versione include un C (C08) e rafforza il pacchetto con A21, restando entro il limite: sforzo totale 3258 + 506 = 3764 (sarebbe troppo) quindi rimuovo A82 per rientrare. Nuova selezione con C incluso e dentro il limite: \{A11, B33, B31, B21, B20, C08, A21\} con sforzo 338+122+231+720+91+131+506 = 2139. Abbiamo margine per aggiungere un altro elemento; propongo aggiungere B09 (552) per salire a 2691, oppure A60 (651) per 2790, restando comunque sotto il limite. Preferisco aggiungere B09.'\}} \\
\texttt{PROPOSTA: \{'A11', 'B33', 'B31', 'B21', 'B20', 'C08', 'A21', 'B09'\}} \\
\texttt{ARGOMENTAZIONE: \{'Sforzo totale 2691: include il C richiesto, mantiene i B ad alto impatto e massimizza l’efficienza senza superare il limite. Se desideri un C ancora più “alto”, possiamo valutare C81 sostituendo due pezzi a basso sforzo, ma questa configurazione è molto solida.'\}} \\
            }
        }
    }
    & \\ \\

    \theutterance \stepcounter{utterance}  
    & & & \multicolumn{4}{p{0.6\linewidth}}{
        \cellcolor[rgb]{0.9,0.9,0.9}{
            \makecell[{{p{\linewidth}}}]{
                \texttt{\tiny{[GM$\langle$P2]}}
                \texttt{RIFIUTO: \{'A11', 'B33', 'B31', 'B21', 'B20', 'C08', 'A21', 'B09'\}} \\
\texttt{ARGOMENTAZIONE: \{'Apprezzo l’inclusione di C08 e il bilanciamento, ma per me è cruciale inserire almeno uno tra C81 o C85. Possiamo farlo rimanendo entro il limite sostituendo alcune B meno prioritarie. Propongo una configurazione che sfrutta al meglio il margine mantenendo buona efficienza complessiva.'\}} \\
\texttt{PROPOSTA: \{'B20', 'B33', 'B31', 'A21', 'A60', 'A82', 'B93', 'C85'\}} \\
\texttt{ARGOMENTAZIONE: \{'Sforzo totale = 91+122+231+506+651+263+492+716 = 3072, ben sotto 3357. Include un C “alto” (C85) e un mix A/B molto produttivo. Se preferisci mantenere B21, posso valutare uno scambio: ad esempio sostituire B93 (492) con B21 (720) e rimuovere A82 (263) per restare sotto il limite (nuovo sforzo 303+? Precisazione: con B21 al posto di B93 e senza A82 si arriva a 303? Correzione: 91+122+231+506+651+720+716 = 3037). Entrambe le varianti sono entro limite; dimmi quale pacchetto B preferisci.'\}} \\
            }
        }
    }
     \\ \\

    \theutterance \stepcounter{utterance}  
    & & \multicolumn{4}{p{0.6\linewidth}}{
        \cellcolor[rgb]{0.9,0.9,0.9}{
            \makecell[{{p{\linewidth}}}]{
                \texttt{\tiny{[P1$\langle$GM]}}
                \texttt{RIFIUTO: \{'A11', 'B33', 'B31', 'B21', 'B20', 'C08', 'A21', 'B09'\}} \\
\texttt{ARGOMENTAZIONE: \{'Apprezzo l’inclusione di C08 e il bilanciamento, ma per me è cruciale inserire almeno uno tra C81 o C85. Possiamo farlo rimanendo entro il limite sostituendo alcune B meno prioritarie. Propongo una configurazione che sfrutta al meglio il margine mantenendo buona efficienza complessiva.'\}} \\
\texttt{PROPOSTA: \{'B20', 'B33', 'B31', 'A21', 'A60', 'A82', 'B93', 'C85'\}} \\
\texttt{ARGOMENTAZIONE: \{'Sforzo totale = 91+122+231+506+651+263+492+716 = 3072, ben sotto 3357. Include un C “alto” (C85) e un mix A/B molto produttivo. Se preferisci mantenere B21, posso valutare uno scambio: ad esempio sostituire B93 (492) con B21 (720) e rimuovere A82 (263) per restare sotto il limite (nuovo sforzo 303+? Precisazione: con B21 al posto di B93 e senza A82 si arriva a 303? Correzione: 91+122+231+506+651+720+716 = 3037). Entrambe le varianti sono entro limite; dimmi quale pacchetto B preferisci.'\}} \\
            }
        }
    }
    & \\ \\

    \theutterance \stepcounter{utterance}  
    & \multicolumn{4}{p{0.6\linewidth}}{
        \cellcolor[rgb]{0.9,0.9,0.9}{
            \makecell[{{p{\linewidth}}}]{
                \texttt{\tiny{[P1$\rangle$GM]}}
                \texttt{RIFIUTO: \{'B20', 'B33', 'B31', 'A21', 'A60', 'A82', 'B93', 'C85'\}} \\
\texttt{ARGOMENTAZIONE: \{'Accolgo l’inclusione di C85, ma questa selezione sacrifica troppi elementi B/A per me cruciali (in particolare A11, B21 e B09) che hanno ottimo rendimento rispetto allo sforzo. Possiamo mantenere un C “alto” e migliorare l’efficienza complessiva sostituendo alcune scelte meno impattanti.'\}} \\
\texttt{PROPOSTA: \{'A11', 'B33', 'B31', 'B21', 'B20', 'C85', 'B09', 'A82'\}} \\
\texttt{ARGOMENTAZIONE: \{'Sforzo totale = 338+122+231+720+91+716+552+263 = 3033, sotto il limite 3357. Include il C alto richiesto (C85) e conserva per me pilastri ad alto impatto (A11, B33, B21, B09) con ottimo rapporto valore/sforzo. Se preferisci C81 al posto di C85, variante possibile: sostituire C85 con C81 e rimuovere A82; nuovo set \{A11, B33, B31, B21, B20, C81, B09\} con sforzo 338+122+231+720+91+849+552 = 2903. Sono aperto a uno scambio puntuale (es. A82 {-}> A21) restando entro il tetto.'\}} \\
            }
        }
    }
    & & \\ \\

    \theutterance \stepcounter{utterance}  
    & & \multicolumn{4}{p{0.6\linewidth}}{
        \cellcolor[rgb]{0.9,0.9,0.9}{
            \makecell[{{p{\linewidth}}}]{
                \texttt{\tiny{[GM$\rangle$P2]}}
                \texttt{RIFIUTO: \{'B20', 'B33', 'B31', 'A21', 'A60', 'A82', 'B93', 'C85'\}} \\
\texttt{ARGOMENTAZIONE: \{'Accolgo l’inclusione di C85, ma questa selezione sacrifica troppi elementi B/A per me cruciali (in particolare A11, B21 e B09) che hanno ottimo rendimento rispetto allo sforzo. Possiamo mantenere un C “alto” e migliorare l’efficienza complessiva sostituendo alcune scelte meno impattanti.'\}} \\
\texttt{PROPOSTA: \{'A11', 'B33', 'B31', 'B21', 'B20', 'C85', 'B09', 'A82'\}} \\
\texttt{ARGOMENTAZIONE: \{'Sforzo totale = 338+122+231+720+91+716+552+263 = 3033, sotto il limite 3357. Include il C alto richiesto (C85) e conserva per me pilastri ad alto impatto (A11, B33, B21, B09) con ottimo rapporto valore/sforzo. Se preferisci C81 al posto di C85, variante possibile: sostituire C85 con C81 e rimuovere A82; nuovo set \{A11, B33, B31, B21, B20, C81, B09\} con sforzo 338+122+231+720+91+849+552 = 2903. Sono aperto a uno scambio puntuale (es. A82 {-}> A21) restando entro il tetto.'\}} \\
            }
        }
    }
    & \\ \\

    \theutterance \stepcounter{utterance}  
    & & & \multicolumn{4}{p{0.6\linewidth}}{
        \cellcolor[rgb]{0.9,0.9,0.9}{
            \makecell[{{p{\linewidth}}}]{
                \texttt{\tiny{[GM$\langle$P2]}}
                \texttt{ARGOMENTAZIONE: \{'Questa proposta è solida e integra il C “alto” che cercavo mantenendo un mix efficiente. Per me sarebbe ideale sostituire A82 con A21: restiamo sotto il limite e incremento leggermente l’allineamento alle mie priorità senza intaccare i tuoi pilastri.'\}} \\
\texttt{PROPOSTA: \{'A11', 'B33', 'B31', 'B21', 'B20', 'C85', 'B09', 'A21'\}} \\
            }
        }
    }
     \\ \\

    \theutterance \stepcounter{utterance}  
    & & \multicolumn{4}{p{0.6\linewidth}}{
        \cellcolor[rgb]{0.9,0.9,0.9}{
            \makecell[{{p{\linewidth}}}]{
                \texttt{\tiny{[P1$\langle$GM]}}
                \texttt{ARGOMENTAZIONE: \{'Questa proposta è solida e integra il C “alto” che cercavo mantenendo un mix efficiente. Per me sarebbe ideale sostituire A82 con A21: restiamo sotto il limite e incremento leggermente l’allineamento alle mie priorità senza intaccare i tuoi pilastri.'\}} \\
\texttt{PROPOSTA: \{'A11', 'B33', 'B31', 'B21', 'B20', 'C85', 'B09', 'A21'\}} \\
            }
        }
    }
    & \\ \\

    \theutterance \stepcounter{utterance}  
    & \multicolumn{4}{p{0.6\linewidth}}{
        \cellcolor[rgb]{0.9,0.9,0.9}{
            \makecell[{{p{\linewidth}}}]{
                \texttt{\tiny{[P1$\rangle$GM]}}
                \texttt{ACCORDO: \{'A11', 'B33', 'B31', 'B21', 'B20', 'C85', 'B09', 'A21'\}} \\
\texttt{ARGOMENTAZIONE: \{'Confermo: sforzo totale = 338+122+231+720+91+716+552+506 = 3276, entro il limite 3357. Mantiene i miei pilastri e include il C “alto” richiesto. Ottimo compromesso, accetto.'\}} \\
            }
        }
    }
    & & \\ \\

    \theutterance \stepcounter{utterance}  
    & & & \multicolumn{2}{p{0.3\linewidth}}{
        \cellcolor[rgb]{0.9,0.9,0.9}{
            \makecell[{{p{\linewidth}}}]{
                \texttt{\tiny{[GM$|$GM]}}
                \texttt{game successful} \\
            }
        }
    }
    & & \\ \\

    \theutterance \stepcounter{utterance}  
    & & & \multicolumn{2}{p{0.3\linewidth}}{
        \cellcolor[rgb]{0.9,0.9,0.9}{
            \makecell[{{p{\linewidth}}}]{
                \texttt{\tiny{[GM$|$GM]}}
                \texttt{end game} \\
            }
        }
    }
    & & \\ \\

\end{supertabular}
}

\end{document}
